% Nouveau EGA V (cote 25): archive cote 25, site page 90 (PDF page 91).
% typescript. Draft compared: none.
% AI-transcribed by Claude Opus 5.5 (claude-opus-5-5, adaptive thinking, effort high), Message Batches API, batch ega25b (msgbatch_01A7n5YuXubUPNK1zsepi8RG), 2026-11-13; not yet checked by eye.
\scanpage{25/90}
\mspage{\unsure{87}}

\uline{Proposition} 14.5. Supposons dim X $\leqslant$ m-1, \var{\xout{Soit Y l'image fermée}}{}{typewriter strike-out}{ega25_25_090-v0001}
\var{\xout{\illegible{1}}}{}{typewriter strike-out}{ega25_25_090-v0002} f:X $\to$ P une immersion, \var{\xout{Soit}}{}{typewriter strike-out}{ega25_25_090-v0003} enfin X séparable sur k. Soit Y\hand{$_y$} l'image
schématique de $X_{k(\hand{y})}$ dans $Q(C_{\hand{y}})$. Alors le morphisme induit
$p'_{C_y} : X_{k(y)} \to Y_y$ est \uline{birationnel}\hand{, et pour tout point x de $X_{k(y)}$ au dessus}
\hand{d'un pt fermé de X, $p'_{C_y}$ est \uline{étale} en x et \unsure{$\overline{m}$|nr} en les points de $p'^{-1}_{C_y}p'_{C_y}(x)$.}
Signalons la conséquence suivante:

\uline{Corollaire} 14.6. Soit X un schéma algébrique projectif \var{\inserted[typed]{irréductible}}{irréductible}{insertion}{ega25_25_090-v0004} \var{\xout{irréductible}}{}{typewriter strike-out}{ega25_25_090-v0005}
\uline{et séparable}
de dimension n sur un corps infini k. Alors \var{\deleted{X est birationellement}}{}{deletion}{ega25_25_090-v0006} \var{\inserted{$\exists$ un plongement birationnel de X}}{$\exists$ un plongement birationnel de X}{insertion}{ega25_25_090-v0007}
\var{\deleted{isomorphe à}}{}{deletion}{ega25_25_090-v0008} \var{\inserted{\var{\deleted{\unsure{dans}}}{}{deletion}{ega25_25_090-v0010} sur}}{\var{\deleted{\unsure{dans}}}{}{deletion}{ega25_25_090-v0010} sur}{insertion}{ega25_25_090-v0009} une hypersurface \var{\deleted{dans}}{}{deletion}{ega25_25_090-v0011} P$^{n+1}$ .

On se gardera de croire \var{\deleted{que}}{}{deletion}{ega25_25_090-v0012}, même si X est un sous-préschéma fermé
lisse géom irréd de P, de dimension m-1 = n , \var{\inserted[typed]{que}}{que}{insertion}{ega25_25_090-v0013} la projection conique p$_{C_{\hand{y}}}$
soit \var{\inserted[typed]{nécessairement}}{nécessairement}{insertion}{ega25_25_090-v0014} une \uline{immersion} . En effet, lorsque k est infini, cela impliquerait
qu'il existe une C rationnelle sur k ayant la même propriété, donc que X
est isomorphe à une hypersurface non singulière de \var{\xout{PxxxQxx}}{}{typewriter strike-out}{ega25_25_090-v0015} P$^{n+1}$. Or
même pour n=1, (\var{\inserted{dans X}}{dans X}{insertion}{ega25_25_090-v0016}courbe algébrique \var{\inserted[typed]{projective}}{projective}{insertion}{ega25_25_090-v0017} lisse \var{\xout{\unsure{sur}}}{}{typewriter strike-out}{ega25_25_090-v0018} et connexe sur un corps alg.
clos), il est facile de construire des exemples où X ne peut s'immerger
dans un P$^2$ . De même, dans 14.4. , on se gardera de confondre l'énon-%
cé donné avec l'assertion (en général fausse) que p$_C$ soit lui-même un
monomorphisme (contre-exemple précédent\var{\xout{\unsure{, xx que}}}{}{typewriter strike-out}{ega25_25_090-v0019}, ou contre-exemple enco-%
re plus évident si X est lisse de dimension m), ou que p$_C$ soit non
ramifié. Pour ce dernier point, on prendra pour s'en convaincre X
sous-\var{\xout{pré}}{}{typewriter strike-out}{ega25_25_090-v0020}schéma fermé lisse \var{\xout{de dim}}{}{typewriter strike-out}{ega25_25_090-v0021} irréd de dim m (sur k alg clos disons
si \var{\xout{XxxxxQx}}{}{typewriter strike-out}{ega25_25_090-v0022} on avait un X $\to$ Q $\simeq$ P$^m$ non ramifié, il serait étale pour
raisons de dimensions, or on peut montrer (cf Chap VIII !) que cela impli-%
que X $\xrightarrow{\sim}$ P$^m$ (P$^m$ étant "simplement connexe). \hand{--} La signification géométri-%
que intuitive de 14.4. est que\var{\xout{x pourxxxvariablexdans}}{}{typewriter strike-out}{ega25_25_090-v0023} l'ensemble de
ramification de p$_{C_{\hand{y}}}$ est ``variable'' \var{\inserted{sur k}}{sur k}{insertion}{ega25_25_090-v0024}, plus précisément l'ensemble de
ramification de p$_{C_{\hand{y}}}$ pour \hand{y} variable dans \var{\inserted{un ouvert de}}{un ouvert de}{insertion}{ega25_25_090-v0025} Grass$^{m+1}$(\hand{$\bar{k}$}) varie dans X(\hand{$\bar{k}$}), et
n'\uline{admet aucun} \var{\deleted{point fermé de X comme}}{}{deletion}{ega25_25_090-v0026} \uline{``point fixe''} ....

\mspage{90}
